\documentclass[12pt]{article}
\usepackage{amsmath, amssymb}
\usepackage[margin=1in]{geometry}
\setlength{\parskip}{6pt}
\title{QUBO Formulation: Exam Timetabling Problem}
\date{}
\begin{document}
\maketitle

\subsection*{Exam Timetabling Problem}

\paragraph{Problem Statement.}
Given a set of $N$ exams indexed by $i \in \{0, \dots, N-1\}$ and a set of $T$ time slots indexed by $t \in \{0, \dots, T-1\}$, the goal is to assign each exam to exactly one time slot. Each pair of exams $(i, j)$ with $i < j$ shares a set of students, quantified by a non-negative conflict weight $w_{ij} \ge 0$. Additionally, each exam $i$ may incur a soft penalty $c_{i,t} \ge 0$ if assigned to slot $t$, representing undesirable scheduling preferences. The objective is to find an assignment that minimizes the total number of student conflicts plus the sum of soft penalties, while strictly satisfying the constraint that every exam is scheduled in exactly one slot.

\paragraph{Decision Variables.}
We introduce binary variables $x_{i,t} \in \{0, 1\}$ for all $i \in \{0, \dots, N-1\}$ and $t \in \{0, \dots, T-1\}$. The variable $x_{i,t}$ equals $1$ if and only if exam $i$ is assigned to time slot $t$, and $0$ otherwise. The total number of variables is $n = N \times T$.

\paragraph{QUBO Derivation.}
We construct a Quadratic Unconstrained Binary Optimization (QUBO) formulation by encoding the objective function and constraints into a single Hamiltonian $H(\mathbf{x})$.

\textbf{Step 1: Objective Function.}
The original objective is to minimize the sum of conflict penalties and soft slot penalties:
\begin{align}
    \min_{\mathbf{x}} \quad H_{\text{obj}}(\mathbf{x}) = \sum_{i=0}^{N-1} \sum_{j=i+1}^{N-1} \sum_{t=0}^{T-1} w_{ij} x_{i,t} x_{j,t} + \sum_{i=0}^{N-1} \sum_{t=0}^{T-1} c_{i,t} x_{i,t}.
\end{align}
Since QUBO minimization is equivalent to the original problem, we retain $H_{\text{obj}}$ as the energy contribution from the objective.

\textbf{Step 2: Constraints.}
Each exam must be assigned to exactly one slot. This is enforced by the equality constraints:
\begin{align}
    \sum_{t=0}^{T-1} x_{i,t} = 1, \quad \forall i \in \{0, \dots, N-1\}.
\end{align}

\textbf{Step 3: Penalty Terms.}
We convert each constraint into a quadratic penalty term using a positive weight $\lambda$. For each exam $i$, the penalty is:
\begin{align}
    H_{\text{pen}, i}(\mathbf{x}) = \lambda \left( \sum_{t=0}^{T-1} x_{i,t} - 1 \right)^2.
\end{align}
Expanding the square and using the idempotent property of binary variables ($x_{i,t}^2 = x_{i,t}$):
\begin{align}
    \left( \sum_{t=0}^{T-1} x_{i,t} - 1 \right)^2 &= \sum_{t=0}^{T-1} x_{i,t}^2 + \sum_{t \neq t'} x_{i,t} x_{i,t'} - 2 \sum_{t=0}^{T-1} x_{i,t} + 1 \nonumber \\
    &= \sum_{t=0}^{T-1} x_{i,t} + \sum_{t \neq t'} x_{i,t} x_{i,t'} - 2 \sum_{t=0}^{T-1} x_{i,t} + 1 \nonumber \\
    &= \sum_{t \neq t'} x_{i,t} x_{i,t'} - \sum_{t=0}^{T-1} x_{i,t} + 1.
\end{align}
Noting that $\sum_{t \neq t'} x_{i,t} x_{i,t'} = 2 \sum_{t < t'} x_{i,t} x_{i,t'}$, the penalty simplifies to:
\begin{align}
    H_{\text{pen}, i}(\mathbf{x}) = \lambda \left( 2 \sum_{0 \le t < t' \le T-1} x_{i,t} x_{i,t'} - \sum_{t=0}^{T-1} x_{i,t} + 1 \right).
\end{align}

\textbf{Step 4: Total Hamiltonian.}
The total QUBO Hamiltonian is the sum of the objective and all penalty terms:
\begin{align}
    H(\mathbf{x}) = H_{\text{obj}}(\mathbf{x}) + \sum_{i=0}^{N-1} H_{\text{pen}, i}(\mathbf{x}).
\end{align}
Substituting the expressions:
\begin{align}
    H(\mathbf{x}) = \sum_{i<j} \sum_{t} w_{ij} x_{i,t} x_{j,t} + \sum_{i,t} c_{i,t} x_{i,t} + \lambda \sum_{i} \left( 2 \sum_{t < t'} x_{i,t} x_{i,t'} - \sum_{t} x_{i,t} + 1 \right).
\end{align}

\textbf{Step 5: Q Matrix and Offset.}
We map each variable $x_{i,t}$ to a unique index $k = i T + t \in \{0, \dots, n-1\}$. The Hamiltonian takes the standard QUBO form $H(\mathbf{x}) = \sum_{k} Q_{k,k} x_k + \sum_{k < k'} Q_{k,k'} x_k x_{k'} + c$. By collecting coefficients from the expanded $H(\mathbf{x})$, we obtain the following closed-form expressions for the Q matrix entries and offset $c$:

\begin{align}
    Q_{iT+t, \, jT+t} &= w_{ij}, \quad \forall i < j, \, \forall t \in \{0, \dots, T-1\}, \\
    Q_{iT+t, \, iT+t'} &= 2\lambda, \quad \forall i \in \{0, \dots, N-1\}, \, \forall t < t', \\
    Q_{iT+t, \, iT+t} &= c_{i,t} - \lambda, \quad \forall i \in \{0, \dots, N-1\}, \, \forall t \in \{0, \dots, T-1\}, \\
    Q_{k, k'} &= 0, \quad \text{otherwise}, \\
    c &= N \lambda.
\end{align}

\paragraph{Penalty Weight Justification.}
The penalty weight $\lambda$ must be chosen sufficiently large to ensure that any violation of the one-slot-per-exam constraint results in an energy increase that cannot be compensated by a reduction in the objective function. A sufficient condition is:
\begin{align}
    \lambda > \max_{i} \left( \sum_{t=0}^{T-1} |c_{i,t}| + \sum_{j \neq i} \sum_{t=0}^{T-1} |w_{ij}| \right).
\end{align}
Under this condition, the minimum penalty incurred by violating the constraint for any exam $i$ (which is at least $\lambda$) strictly exceeds the maximum possible reduction in the objective terms involving exam $i$. Consequently, the global minimum of $H(\mathbf{x})$ must correspond to a feasible assignment.

\paragraph{Correctness Argument.}
Minimizing $H(\mathbf{x})$ over $\{0,1\}^n$ recovers the optimal exam schedule as follows. For any feasible assignment, the constraint terms vanish, and $H(\mathbf{x}) = H_{\text{obj}}(\mathbf{x})$. For any infeasible assignment, at least one constraint is violated, incurring a penalty of at least $\lambda$. By the choice of $\lambda$, the energy of any infeasible solution exceeds that of the best feasible solution. Therefore, the global minimum of $H(\mathbf{x})$ is attained at a feasible point that minimizes $H_{\text{obj}}(\mathbf{x})$, which is precisely the optimal solution to the original timetabling problem.

\paragraph{Illustrative Example.}
Consider a small instance with $N=2$ exams and $T=2$ slots. Let the conflict weight be $w_{0,1} = 1$, and all soft penalties be zero ($c_{i,t} = 0$). The variables are $x_{0,0}, x_{0,1}, x_{1,0}, x_{1,1}$, mapped to indices $0, 1, 2, 3$ respectively.

The objective terms are:
\begin{align}
    H_{\text{obj}} = 1 \cdot x_{0,0} x_{1,0} + 1 \cdot x_{0,1} x_{1,1} = x_0 x_2 + x_1 x_3.
\end{align}
The penalty terms for $\lambda = 1$ are:
\begin{align}
    H_{\text{pen}, 0} &= 1 \cdot (2 x_0 x_1 - x_0 - x_1 + 1), \\
    H_{\text{pen}, 1} &= 1 \cdot (2 x_2 x_3 - x_2 - x_3 + 1).
\end{align}
The total Hamiltonian is:
\begin{align}
    H(\mathbf{x}) = x_0 x_2 + x_1 x_3 + 2 x_0 x_1 - x_0 - x_1 + 2 x_2 x_3 - x_2 - x_3 + 2.
\end{align}
The Q matrix entries are $Q_{0,2}=1, Q_{1,3}=1, Q_{0,1}=2, Q_{2,3}=2$, and diagonal entries $Q_{0,0}=Q_{1,1}=Q_{2,2}=Q_{3,3}=-1$. The offset is $c=2$.

Evaluating feasible solutions:
\begin{itemize}
    \item $x = (1,0,0,1)$ (exams in slots 0 and 1): $H = 0 + 0 + 0 - 1 - 0 + 0 - 0 - 1 + 2 = 0$.
    \item $x = (0,1,1,0)$ (exams in slots 1 and 0): $H = 0 + 0 + 0 - 0 - 1 + 0 - 1 - 0 + 2 = 0$.
    \item $x = (1,0,1,0)$ (both in slot 0): $H = 1 + 0 + 0 - 1 - 0 + 0 - 1 - 0 + 2 = 1$.
\end{itemize}
The minimum energy is $0$, achieved by assigning the exams to different slots, which correctly resolves the conflict.
\end{document}